\documentclass[11pt]{article}

\usepackage[margin=1in]{geometry}
\usepackage{booktabs}
\usepackage{hyperref}
\usepackage{amsmath,amssymb}
\usepackage{enumitem}

\newcommand{\Etheta}{E_\theta}
\newcommand{\Ztheta}{Z_\theta}

\title{Free-Energy Inference for Verifier-Guided Structured Generation}
\author{Brando Miranda}
\date{\today \\ \emph{Planning draft}}

\begin{document}
\maketitle

\begin{abstract}
This paper is the conditional upside paper for the Free Energy project.
It should be advanced only if the proposed novel EBM produces a concrete
advantage over autoregressive and conventional EBM baselines on verifier-guided
structured generation.
\end{abstract}

\section{Problem}

Autoregressive models turn global structured generation into local next-token
prediction.
Energy-based models score whole objects but inherit hard normalization and
sampling.
The goal of this paper is to propose and test a model that keeps the global
compatibility benefits of energies while making training and inference practical
enough for verifier-guided tasks.

\section{Model}

The model should be stated as an energy or free-energy objective.
At minimum, specify:

\begin{equation}
  \Etheta(x, y)
\end{equation}

for context $x$ and candidate structured object $y$, plus the training surrogate
that avoids or approximates

\begin{equation}
  \Ztheta(x) = \sum_y \exp(-\Etheta(x,y)).
\end{equation}

\section{Training objective}

The paper should state exactly which cost is being avoided or relocated:

\begin{itemize}[leftmargin=2em]
  \item local softmax normalization;
  \item global partition-function estimation;
  \item verifier-guided contrastive or margin loss;
  \item free-energy/variational surrogate;
  \item test-time energy descent or search.
\end{itemize}

\section{Inference}

Inference should be described as a budgeted optimization procedure rather than a
vague global search.
Report the number of energy evaluations, verifier calls, gradient steps, samples,
and retries.

\section{Experiments}

\begin{enumerate}[leftmargin=2em]
  \item Toy controls where the mechanism is known to help.
  \item VeriBench/Lean structured generation.
  \item MNIST/order-effect smoke if it clarifies the mechanism.
  \item Ablations against AR and conventional EBM baselines.
\end{enumerate}

\section{Go/no-go criteria}

This paper should not be submitted unless at least one of the following is true:

\begin{itemize}[leftmargin=2em]
  \item pass@k improves over both AR and normal EBM baselines at matched budget;
  \item cost per verified solution improves substantially;
  \item scaling with proof/program difficulty is better;
  \item a negative result reveals a publishable failure mode of existing EBM
        assumptions.
\end{itemize}

\end{document}
